\section{Task 3: Algorithm Design --- SSD-R}

\textbf{SSD-R: SLA-Slack DSATUR with Certified infeasibility and Repair.}
SSD-R is a hybrid with four stages. Stage 1 certifies infeasibility using
necessary conditions. Stage 2 is a saturation-ordered constructive greedy
(DSATUR generalised to windows and capacity) with priority restarts.
Stage 3 is a one-level eviction repair. Stage 4 is a local search on $P$
using moves and swaps. Labels L$x$ and R$x$ below appear as comments in
\texttt{solver/algorithm.py}, so the code can be matched to the pseudocode
line by line.

\subsection{Pseudocode}
Notation: $N(i)$ is the set of conflict neighbours of $t_i$, and
$\text{load}(s)\in\mathbb{R}^d$ is the current load of slot $s$.
$F(i)=\{s\in[l_i,u_i] : N(i)\cap\sigma^{-1}(s)=\emptyset,\ \text{load}(s)+r_i\le C(s)\}$
is the set of slots that are currently feasible for $t_i$.
$c(i,s) = w_i(s+1) + \lambda_R w_i\big(\tfrac{s-l_i}{u_i-l_i}\big)^2$ is the
separable part of $P$, and $\Delta B$ is the change in the imbalance term.

\begin{algorithm}[h]
\caption{SSD-R$(T, G, r, C, w, \tau)$}
\begin{algorithmic}
\State \textbf{Stage 1 --- certify infeasibility}
\State L1:\ \textbf{for} each $t_i$: \textbf{if} no $s\in[l_i,u_i]$ has $r_i\le C(s)$ \textbf{return} INFEASIBLE \Comment{fits nowhere}
\State L2:\ \textbf{for} each interval $[a,b]\subseteq[0,K-1]$: $S\gets\{i: a\le l_i,\ u_i\le b\}$
\State L3:\ \quad \textbf{if} $\exists k: \sum_{i\in S} r_{ik} > \sum_{s=a}^{b} C_k(s)$ \textbf{return} INFEASIBLE \Comment{interval demand}
\State L4:\ \quad \textbf{if} $G[S]$ is visibly not $(b{-}a{+}1)$-colourable \textbf{return} INFEASIBLE
\Statex \hspace{2.6em}\small(width 1: any edge; width 2: odd cycle by BFS; width $\ge3$: greedy clique $>b{-}a{+}1$)\normalsize
\State \textbf{Stage 2 --- construct, with restarts}
\State R1:\ $\Pi\gets[\,]$; \textbf{repeat} at most $n$ times: \Comment{$\Pi$: promoted tasks}
\State R2:\ \quad run L5--L11 with promotion list $\Pi$; \textbf{if} it succeeds \textbf{goto} L18
\State R3:\ \quad move the failed task to the front of $\Pi$
\State \quad \textbf{return} NOT\_FOUND (best partial $\sigma$)
\State L5:\ \quad $U\gets T$, $\sigma\gets\emptyset$
\State L6:\ \quad \textbf{while} $U\neq\emptyset$:
\State L7:\ \qquad compute $F(i)$ for all $i\in U$
\State L8:\ \qquad $i^\ast\gets\arg\min_{i\in U}\big(\text{rank}_\Pi(i),\ |F(i)|,\ -\deg(i),\ -w_i,\ i\big)$
\State L9:\ \qquad \textbf{if} $F(i^\ast)=\emptyset$: $F(i^\ast)\gets\textsc{Repair}(i^\ast)$; \textbf{if} still $\emptyset$ \textbf{fail with} $i^\ast$
\State L10:\ \qquad $S\gets\{s\in F(i^\ast): \nexists j\in U{\setminus}\{i^\ast\}$ with $F(j)=\{s\}$ that $i^\ast$ would block$\}$, or $F(i^\ast)$ if empty
\Statex \hspace{3.6em} $\sigma(i^\ast)\gets\arg\min_{s\in S}\big(c(i^\ast,s)+\lambda_B\Delta B,\ s\big)$ \Comment{forward check, then cheapest}
\State L11:\ \qquad $U\gets U\setminus\{i^\ast\}$
\State \textbf{Stage 3 --- }\textsc{Repair}$(i)$
\State L12:\ \textbf{for} $s=l_i,\dots,u_i$:
\State L13:\ \quad $B\gets N(i)\cap\sigma^{-1}(s)$ \Comment{conflict blockers}
\State L14:\ \quad \textbf{while} $r_i$ does not fit in $s$ without $B$: add to $B$ the largest task of $s$ \Comment{overflowing dim}
\State L15:\ \quad tentatively move each $j\in B$ to its cheapest feasible slot $\neq s$
\State L16:\ \quad \textbf{if} all moved and $s\in F(i)$ \textbf{return} $\{s\}$ \textbf{else} roll back
\State L17:\ \textbf{return} $\emptyset$
\State \textbf{Stage 4 --- improve}
\State L18:\ \textbf{repeat} until no improvement (at most 50 passes):
\State L19:\ \quad for each $i$ and $s\in[l_i,u_i]$: move $i\to s$ if feasible and $\Delta P<-\varepsilon$
\State L20:\ \quad for each pair $i,j$ with $\sigma(i)\ne\sigma(j)$: swap if feasible and $\Delta P<-\varepsilon$
\State L21:\ \textbf{return} $\sigma$
\end{algorithmic}
\end{algorithm}

\subsection{Line-level justification}
\begin{description}
\item[L1, L3] Each is a necessary condition. If one fails, the failure is a
  checkable proof of infeasibility (Lemma~\ref{lem:cert}). L3 covers every
  interval, not only $[0,K-1]$. Tasks with narrow windows can overload a
  sub-range even when the global totals fit.
\item[L4] Tasks confined to $[a,b]$ must be coloured with $b-a+1$ colours.
  For width 2 this is exactly bipartiteness, which BFS decides in linear
  time, so the check is \emph{exact}. Across the 9 benchmarks, 31\% of
  windows have width 2, and 50\% have width at most 3, so this case
  matters most in practice. For width $\ge 3$,
  $k$-colourability is NP-hard. A clique larger than the width is a sound
  but incomplete certificate; it is grown greedily from every start vertex.
\item[L8] This is the core DSATUR idea, adapted to this problem. A task's
  ``saturation'' is the number of slots \emph{still} open to it, which
  shrinks through conflicts (F1), capacity (F2) \emph{and} the window (F3).
  The most constrained task is placed first, while it still has options.
  Ties go to higher degree, because that task constrains the most others,
  then to heavier weight, because it should get the earlier slot. Index
  breaks the final tie, so runs are deterministic.
\item[L10] First a forward check: a slot is skipped if it is the
  \emph{last} feasible slot of another unplaced task $j$, and $i^\ast$
  would block it there (a conflict with $j$, or no capacity left for both).
  Among the remaining slots, the task takes the one that adds the least to
  $P$. Without the forward check, the cheapest-slot rule repeatedly walked
  into dead ends. On $(100,10,0.05)$ seed 1, four tasks with nested windows
  evicted each other in a cycle through all $n$ restarts (\S6). The tie
  rule ``lower slot'' keeps zero-weight tasks deterministic (used in the
  tight example, \S4).
\item[L14] If capacity, not a conflict, is what blocks $i$, the largest
  task in the overflowing dimension is evicted first, so that as few tasks
  as possible move.
\item[L15--L16] Repair goes only one level deep: evicted tasks must go to
  slots that are already feasible for them. Chains of evictions would turn
  repair into a full search. All moves are rolled back on failure, so the
  partial assignment stays feasible at all times.
\item[R1--R3] Squeaky-wheel restarts. A task that fails late was usually
  squeezed out by earlier choices, so it is promoted to go first in the
  next pass. On 1200 random small instances, restarts cut the feasible
  instances that SSD-R missed from 58 to 2 (\S6).
\item[L19--L20] A move lowers $P$ directly. A swap fixes cases that no
  single move can, for example when a full slot holds a light task and a
  heavy task waits in a later slot. Only strict improvements are accepted,
  so $P$ falls at every step and the search cannot cycle.
\end{description}

\subsection{Complexity}
Stage 1 checks $O(K^2)$ intervals. Each costs $O(nd)$ for demand and
$O(n^2\omega)$ for the clique, where $\omega$ is the clique size. One pass
of L5--L11 does $n$ iterations, each computing $F$ for at most $n$ tasks in
$O(Kd)$ using a per-(task, slot) conflict counter. Repair costs
$O(K\cdot n\cdot Kd)$. With at most $n$ restarts, Stage 2 is
$O(n^3Kd + n^2K^2d)$. One local-search pass is $O(n^2 d)$, with at most 50
passes. Every term is polynomial in $n, K, d$. Measured runtimes are
in~\S6.

\subsection{Design rationale}
The problem is DSATUR-shaped because, on generator instances, the binding
constraint is almost always the \emph{interaction of narrow SLA windows
with conflicts}. Most windows span only 2--3 slots, so a task with a few
neighbours already in its window has no slot left. A greedy that ignores
how many options each task has left (for example, heaviest weight first)
places flexible tasks first and strands the constrained ones. Saturation
ordering places each task while it still has a choice.

Stage 1 comes first because infeasibility is \emph{common} on these
instances. Six of the nine graded benchmarks are infeasible (\S6). A
heuristic that only reports ``not found'' cannot tell a hard instance from
an impossible one. Certificates make that distinction and give the
operator an actionable reason, such as ``these 5 bureau pulls form an odd
conflict cycle inside a 2-slot SLA window''.

Penalty-awareness appears in two places. L10 is a myopic cost-aware
choice. Stage 4 recovers the global trade-offs, such as delay against SLA
risk against load balance, that a one-pass greedy cannot see.

\subsection{Alternatives considered and rejected}
\begin{enumerate}
\item \textbf{Exact branch-and-bound (or an ILP).} This gives the optimal
  $\sigma$, and the infeasibility proof is exhaustive. We implemented it
  (\texttt{solver/bruteforce.py}) and use it as ground truth. It was
  rejected as the main algorithm because it is exponential: $K^n$ leaves in
  the worst case, which is $20^{200}$ for the sparse stress instance.
  Writing it as an ILP would also need a forbidden solver.
\item \textbf{Simulated annealing on a relaxed objective} $P(\sigma) + M
  \cdot \#\text{violations}$. It is easy to write, and the penalty handles
  all three constraint families the same way. It was rejected for three
  reasons. (i) It can never \emph{prove} infeasibility. On the 6 infeasible
  benchmarks it would run to its time limit and report ``best effort, still
  violating'', with no reason. (ii) It is randomised, so a viva trace or a
  production incident cannot be reproduced without the seed, and its
  quality depends on the cooling schedule, $M$ and the iteration budget.
  (iii) It has no worst-case structure from which an approximation ratio
  can be derived.
\end{enumerate}
